% ╔══════════════════════════════════════════════════════════════════╗
% ║  Графовое тестирование: автогенерация тестов из описания системы ║
% ║  Компилировать: xelatex graph_testing_research.tex               ║
% ╚══════════════════════════════════════════════════════════════════╝

\documentclass[12pt, a4paper, twoside]{article}

% ─── Шрифты и язык ───────────────────────────────────────────────
\usepackage{fontspec}
\usepackage{polyglossia}
\setdefaultlanguage{russian}
\setotherlanguage{english}
\setmainfont{Liberation Serif}
\setsansfont{Liberation Sans}
\setmonofont[Scale=0.85]{DejaVu Sans Mono}

% ─── Геометрия страницы ──────────────────────────────────────────
\usepackage[
  a4paper,
  top=25mm,
  bottom=25mm,
  left=30mm,
  right=25mm,
  headheight=25pt
]{geometry}

% ─── Базовые пакеты ──────────────────────────────────────────────
\usepackage[protrusion=false,expansion=false]{microtype}
\usepackage{setspace}
\usepackage{parskip}
\usepackage{booktabs}
\usepackage{longtable}
\usepackage{array}
\usepackage{enumitem}
\usepackage{graphicx}
\usepackage{float}
\usepackage{caption}
\usepackage{subcaption}
\usepackage{amsmath}
\usepackage{amssymb}

% ─── TikZ для графов ─────────────────────────────────────────────
\usepackage{tikz}
\usetikzlibrary{arrows.meta, positioning, fit, backgrounds, shadows.blur, calc, shapes.geometric, automata}

% ─── Цвета ───────────────────────────────────────────────────────
\usepackage[dvipsnames, svgnames, x11names]{xcolor}
\definecolor{AccentBlue}{RGB}{25, 90, 160}
\definecolor{AccentTeal}{RGB}{0, 120, 120}
\definecolor{AccentOrange}{RGB}{200, 80, 0}
\definecolor{LightGray}{RGB}{248, 248, 250}
\definecolor{MidGray}{RGB}{130, 130, 140}
\definecolor{DarkGray}{RGB}{40, 40, 50}
\definecolor{CodeBg}{RGB}{245, 246, 250}
\definecolor{CodeBorder}{RGB}{200, 205, 220}
\definecolor{RuleColor}{RGB}{25, 90, 160}
\definecolor{LinkColor}{RGB}{10, 80, 160}
\definecolor{GreenOk}{RGB}{30, 140, 60}
\definecolor{StateInit}{RGB}{70, 130, 180}
\definecolor{StateTarget}{RGB}{34, 139, 34}
\definecolor{StateForbid}{RGB}{178, 34, 34}
\definecolor{StateNormal}{RGB}{100, 100, 120}

% ─── Гиперссылки ─────────────────────────────────────────────────
\usepackage[
  unicode,
  colorlinks=true,
  linkcolor=LinkColor,
  urlcolor=AccentTeal,
  citecolor=AccentOrange,
  pdftitle={Графовое тестирование: автогенерация тестов из описания системы},
  pdfauthor={Opus Research},
  pdfsubject={development\_system :: Model-Based Testing :: Common Lisp},
  pdfkeywords={Common Lisp, MBT, property-based testing, граф состояний, автотесты}
]{hyperref}

% ─── Код (listings) ──────────────────────────────────────────────
\usepackage{listings}
\usepackage{lstfiracode}
\lstset{
  basicstyle=\small\ttfamily\color{DarkGray},
  backgroundcolor=\color{CodeBg},
  frame=single,
  framerule=0.6pt,
  rulecolor=\color{CodeBorder},
  breaklines=true,
  breakatwhitespace=false,
  showspaces=false,
  showstringspaces=false,
  showtabs=false,
  tabsize=2,
  captionpos=b,
  numbers=left,
  numberstyle=\tiny\color{MidGray},
  numbersep=8pt,
  xleftmargin=14pt,
  xrightmargin=4pt,
  aboveskip=10pt,
  belowskip=6pt,
  escapeinside={(*@}{@*)},
  literate=
    {а}{{\cyra}}1  {б}{{\cyrb}}1  {в}{{\cyrv}}1  {г}{{\cyrg}}1
    {д}{{\cyrd}}1  {е}{{\cyre}}1  {ё}{{\cyryo}}1 {ж}{{\cyrzh}}1
    {з}{{\cyrz}}1  {и}{{\cyri}}1  {й}{{\cyrishrt}}1 {к}{{\cyrk}}1
    {л}{{\cyrl}}1  {м}{{\cyrm}}1  {н}{{\cyrn}}1  {о}{{\cyro}}1
    {п}{{\cyrp}}1  {р}{{\cyrr}}1  {с}{{\cyrs}}1  {т}{{\cyrt}}1
    {у}{{\cyru}}1  {ф}{{\cyrf}}1  {х}{{\cyrh}}1  {ц}{{\cyrc}}1
    {ч}{{\cyrch}}1 {ш}{{\cyrsh}}1 {щ}{{\cyrshch}}1 {ъ}{{\cyrhrdsn}}1
    {ы}{{\cyrery}}1 {ь}{{\cyrsftsn}}1 {э}{{\cyrerev}}1 {ю}{{\cyryu}}1
    {я}{{\cyrya}}1
    {А}{{\CYRA}}1  {Б}{{\CYRB}}1  {В}{{\CYRV}}1  {Г}{{\CYRG}}1
    {Д}{{\CYRD}}1  {Е}{{\CYRE}}1  {Ё}{{\CYRYO}}1 {Ж}{{\CYRZH}}1
    {З}{{\CYRZ}}1  {И}{{\CYRI}}1  {Й}{{\CYRISHRT}}1 {К}{{\CYRK}}1
    {Л}{{\CYRL}}1  {М}{{\CYRM}}1  {Н}{{\CYRN}}1  {О}{{\CYRO}}1
    {П}{{\CYRP}}1  {Р}{{\CYRR}}1  {С}{{\CYRS}}1  {Т}{{\CYRT}}1
    {У}{{\CYRU}}1  {Ф}{{\CYRF}}1  {Х}{{\CYRH}}1  {Ц}{{\CYRC}}1
    {Ч}{{\CYRCH}}1 {Ш}{{\CYRSH}}1 {Щ}{{\CYRSHCH}}1 {Ъ}{{\CYRHRDSN}}1
    {Ы}{{\CYRERY}}1 {Ь}{{\CYRSFTSN}}1 {Э}{{\CYREREV}}1 {Ю}{{\CYRYU}}1
    {Я}{{\CYRYA}}1,
}

\lstdefinelanguage{CommonLisp}{
  morekeywords={defun,defpackage,defparameter,defmacro,defstruct,let,let*,cond,when,unless,
                if,and,or,not,progn,dolist,mapcar,mapcan,reduce,sort,
                case,format,in-package,use,export,import,load,compile-file,
                handler-case,error,return-from,values,multiple-value-bind,
                setf,push,pop,car,cdr,caar,cdar,cons,list,append,remove-if,
                remove-if-not,find,find-if,count-if,subseq,length,first,last,
                second,third,loop,do,dotimes,with-open-file,ensure-directories-exist,
                assert,check-type,the,declare,type,ftype,lambda,funcall,apply,
                gethash,make-hash-table,maphash,remhash,hash-table-count,
                defgeneric,defmethod,defclass,make-instance,slot-value,with-slots,
                for-all,is,signals,finishes,test,def-suite,in-suite,run},
  morestring=[b]",
  morecomment=[l]{;;},
  sensitive=false,
  keywordstyle=\color{AccentBlue}\bfseries,
  commentstyle=\color{MidGray}\itshape,
  stringstyle=\color{AccentOrange},
}

\lstdefinelanguage{JSON}{
  morestring=[b]",
  morestring=[d]',
  stringstyle=\color{AccentOrange},
  literate=
    *{0}{{{\color{AccentBlue}0}}}{1}
     {1}{{{\color{AccentBlue}1}}}{1}
     {2}{{{\color{AccentBlue}2}}}{1}
     {3}{{{\color{AccentBlue}3}}}{1}
     {4}{{{\color{AccentBlue}4}}}{1}
     {5}{{{\color{AccentBlue}5}}}{1}
     {6}{{{\color{AccentBlue}6}}}{1}
     {7}{{{\color{AccentBlue}7}}}{1}
     {8}{{{\color{AccentBlue}8}}}{1}
     {9}{{{\color{AccentBlue}9}}}{1}
     {:}{{{\color{DarkGray}{:}}}}{1}
     {,}{{{\color{DarkGray}{,}}}}{1}
     {\{}{{{\color{AccentBlue}{\{}}}}{1}
     {\}}{{{\color{AccentBlue}{\}}}}}{1}
     {[}{{{\color{AccentTeal}[}}}{1}
     {]}{{{\color{AccentTeal}]}}}{1},
}

% ─── Колонтитулы ─────────────────────────────────────────────────
\usepackage{fancyhdr}
\pagestyle{fancy}
\fancyhf{}
\fancyhead[LE]{\small\color{MidGray}\textit{Графовое тестирование}}
\fancyhead[RO]{\small\color{MidGray}\textit{development\_system}}
\fancyfoot[C]{\small\color{MidGray}\thepage}
\renewcommand{\headrulewidth}{0.4pt}
\renewcommand{\headrule}{\hbox to\headwidth{\color{RuleColor}\leaders\hrule height \headrulewidth\hfill}}
\renewcommand{\footrulewidth}{0pt}

% ─── Заголовки разделов ──────────────────────────────────────────
\usepackage{titlesec}
\titleformat{\section}
  {\large\bfseries\color{AccentBlue}}
  {\thesection.}{0.8em}{}
  [\vspace{2pt}\color{RuleColor}\rule{\linewidth}{0.6pt}\vspace{4pt}]

\titleformat{\subsection}
  {\normalsize\bfseries\color{DarkGray}}
  {\thesubsection.}{0.7em}{}

\titleformat{\subsubsection}
  {\small\bfseries\color{AccentTeal}}
  {\thesubsubsection.}{0.6em}{}

\titlespacing*{\section}{0pt}{18pt}{8pt}
\titlespacing*{\subsection}{0pt}{12pt}{4pt}
\titlespacing*{\subsubsection}{0pt}{8pt}{2pt}

% ─── tcolorbox для callout-блоков ────────────────────────────────
\usepackage{tcolorbox}
\tcbuselibrary{breakable, skins}

\newtcolorbox{infobox}[1][]{
  colback=LightGray,
  colframe=AccentBlue,
  colbacktitle=AccentBlue,
  coltitle=white,
  fonttitle=\small\bfseries,
  boxrule=0.6pt,
  arc=3pt,
  breakable,
  #1
}

\newtcolorbox{warningbox}[1][]{
  colback=Ivory,
  colframe=AccentOrange,
  colbacktitle=AccentOrange,
  coltitle=white,
  fonttitle=\small\bfseries,
  boxrule=0.6pt,
  arc=3pt,
  breakable,
  #1
}

\newtcolorbox{successbox}[1][]{
  colback=Honeydew,
  colframe=GreenOk,
  colbacktitle=GreenOk,
  coltitle=white,
  fonttitle=\small\bfseries,
  boxrule=0.6pt,
  arc=3pt,
  breakable,
  #1
}

% ─── Таблицы ─────────────────────────────────────────────────────
\usepackage{tabularx}
\newcolumntype{L}{>{\raggedright\arraybackslash}X}
\newcolumntype{C}{>{\centering\arraybackslash}X}
\renewcommand{\arraystretch}{1.35}

% ─── Заголовок документа ─────────────────────────────────────────
\usepackage{titling}

% ─── Прочее ──────────────────────────────────────────────────────
\usepackage{soul}
\usepackage{csquotes}
\usepackage{footmisc}
\setlength{\parskip}{6pt}
\setlength{\parindent}{0pt}
\onehalfspacing

% ═══════════════════════════════════════════════════════════════════
\begin{document}

% ─── Титульная страница ──────────────────────────────────────────
\begin{titlepage}
  \newgeometry{margin=30mm}
  \color{DarkGray}

  \vspace*{10mm}

  % Верхняя линия
  {\color{AccentBlue}\rule{\textwidth}{1.5pt}}
  \vspace{6mm}

  % Метка
  {\sffamily\small\color{MidGray}{ТЕХНИЧЕСКОЕ ИССЛЕДОВАНИЕ}}\\[4mm]

  % Главный заголовок
  {\sffamily\LARGE\bfseries\color{AccentBlue}
    Графовое тестирование:\\[4pt]
    автогенерация тестов из описания системы\\[4pt]
    в \texttt{development\_system}
  }

  \vspace{8mm}
  {\color{AccentBlue}\rule{\textwidth}{0.4pt}}
  \vspace{8mm}

  % Аннотация на титульной
  \begin{minipage}{0.92\textwidth}
    \itshape\color{DarkGray}
    Данное исследование описывает методологию автоматической генерации
    тестов на основе графа состояний системы. Человек описывает систему
    как граф (текстом, JSON, или Lisp-DSL), а специализированный поток
    \textbf{граф-тестировщик} автоматически порождает тест-кейсы:
    для целевых состояний --- тесты достижимости, для нежелательных ---
    тесты на исключения. Исследование охватывает теоретические основы
    (Model-Based Testing, Property-Based Testing, Reachability Analysis),
    конкретную архитектуру в рамках существующей метациркулярной системы,
    примеры кода на Common Lisp и пошаговый план реализации.
  \end{minipage}

  \vspace{12mm}

  % Мета-таблица
  \begin{tabular}{ll}
    \textbf{Версия:}  & 1.0 \\[2pt]
    \textbf{Дата:}    & 14 марта 2026 \\[2pt]
    \textbf{Проект:}  & \texttt{development\_system} \\[2pt]
    \textbf{Стек:}    & Common Lisp (SBCL), OpenRouter, Telegram \\[2pt]
    \textbf{Предназначение:} & автоматизация тестирования метациркулярной среды \\
  \end{tabular}

  \vfill

  % Нижняя линия + ключевые слова
  {\color{AccentBlue}\rule{\textwidth}{0.4pt}}
  \vspace{4mm}
  {\small\color{MidGray}
    \textbf{Ключевые слова:}
    Common Lisp,\enspace Model-Based Testing,\enspace Property-Based Testing,\enspace
    Граф состояний,\enspace Reachability,\enspace TLA+,\enspace Fiveam,\enspace
    Метациркулярность,\enspace Автогенерация тестов
  }

  \restoregeometry
\end{titlepage}

% ─── Оглавление ──────────────────────────────────────────────────
\tableofcontents
\newpage

% ═══════════════════════════════════════════════════════════════════
\section{Введение}

\subsection{Постановка задачи}

Артемий формулирует задачу:

\begin{quote}
\itshape
<<Просмотреть граф как система сейчас реализована (описанный человеком) и
реализовать на основе графа автотесты по системе: входные данные такие-то,
целевое состояние должно быть таким-то, оно должно быть достижимо,
а нежелательные состояния должны выбрасывать исключения.>>
\end{quote}

\textbf{Суть задачи} в трёх шагах:

\begin{enumerate}[leftmargin=1.5em]
  \item Человек описывает систему как \textbf{граф состояний/переходов}
        (текстом, схемой, или формально)
  \item Система автоматически \textbf{порождает тесты} вида:
  \begin{itemize}
    \item <<При входе $X$ $\to$ целевое состояние $Y$ \textbf{достижимо}>>
    \item <<При входе $X'$ $\to$ должно выброситься исключение (нежелательное состояние)>>
  \end{itemize}
  \item Тесты \textbf{запускаются} и проверяют реальную систему
\end{enumerate}

\subsection{Контекст системы \texttt{development\_system}}

\texttt{development\_system} --- метациркулярная среда вычислительных потоков:

\begin{center}
\begin{tabular}{ll}
  \toprule
  \textbf{Компонент} & \textbf{Роль} \\
  \midrule
  Telegram         & Входная точка: пользователь отправляет команды \\
  Мастер (SBCL)    & Polling loop, маршрутизация команд, хранение истории \\
  \texttt{команды.lisp} & Обработчики команд (\texttt{/запустить}, \texttt{/обучить}, ...) \\
  \texttt{потоки.lisp}  & Загрузка и запуск потоков через reflection \\
  \texttt{мастер/потоки/*.lisp} & Каждый поток: один пакет, одна функция \texttt{выполнить} \\
  Порождатель      & Метапоток: по описанию генерирует новый поток через LLM \\
  \bottomrule
\end{tabular}
\end{center}

\textbf{Ключевые потоки:}

\begin{itemize}[leftmargin=1.5em]
  \item \texttt{эхо} --- возвращает входную задачу (для отладки)
  \item \texttt{уточнение} --- уточняет задачу через LLM
  \item \texttt{исполнитель} --- выполняет задачу через LLM
  \item \texttt{питон} --- LLM генерирует Python-код, запускает, возвращает вывод
  \item \texttt{порождатель} --- принимает описание на русском, генерирует новый поток
\end{itemize}

\subsection{Механизм загрузки потоков}

\begin{lstlisting}[language=CommonLisp,caption={Reflection-based поиск функции выполнить}]
(defun %найти-выполнить (имя)
  "Находит функцию ВЫПОЛНИТЬ в пакете :ПОТОК-ИМЯ"
  (let* ((pkg (find-package
               (string-upcase (concatenate 'string "поток-" имя))))
         (sym (when pkg (find-symbol "ВЫПОЛНИТЬ" pkg))))
    (when (and sym (fboundp sym))
      (symbol-function sym))))

;; Состояние системы
(defparameter *активные-потоки* (make-hash-table :test #'equal))
\end{lstlisting}

\textbf{Метациркулярность:} порождатель --- сам поток, который генерирует
другие потоки через \texttt{compile-file + load} в живой образ SBCL.

\subsection{Цель исследования}

Спроектировать и реализовать \texttt{поток-граф-тестировщик}, который:

\begin{enumerate}[leftmargin=1.5em]
  \item Принимает описание системы (естественный язык, JSON, Lisp-DSL)
  \item Строит внутреннее представление графа состояний
  \item Автоматически генерирует тест-кейсы через анализ графа
  \item Выполняет тесты через существующую инфраструктуру
  \item Возвращает отчёт: passed / failed / unreachable
\end{enumerate}

% ═══════════════════════════════════════════════════════════════════
\section{Теоретические основы}

\subsection{Model-Based Testing (MBT)}

\textbf{Model-Based Testing} --- методология тестирования, при которой
тест-кейсы автоматически генерируются из формальной модели системы.

\begin{infobox}[title={Ключевая идея MBT}]
Вместо ручного написания тестов инженер создаёт \textbf{модель} системы
(граф состояний, конечный автомат, спецификацию). Алгоритм обходит модель
и генерирует последовательности входов, покрывающие все переходы.
\end{infobox}

\subsubsection{Transition Tree Method}

Метод дерева переходов --- базовый алгоритм MBT:

\begin{enumerate}[leftmargin=1.5em]
  \item Начать с начального состояния (корень дерева)
  \item Для каждого исходящего перехода создать ветку
  \item Рекурсивно повторить для всех достижимых состояний
  \item Каждый путь от корня до листа --- один тест-кейс
\end{enumerate}

\subsubsection{Критерии покрытия}

\begin{center}
\begin{tabularx}{\textwidth}{lLl}
  \toprule
  \textbf{Критерий} & \textbf{Описание} & \textbf{Полнота} \\
  \midrule
  All-states     & Каждое состояние посещается хотя бы раз & Низкая \\
  All-transitions & Каждый переход (ребро) выполняется хотя бы раз & Средняя \\
  All-paths      & Все возможные пути от начала до конца & Высокая \\
  \bottomrule
\end{tabularx}
\end{center}

\textit{Для \texttt{development\_system} рекомендуется:} \textbf{All-transitions} ---
оптимальный баланс между полнотой и трудоёмкостью.

\subsubsection{Negative Testing через граф}

Negative testing --- проверка того, что система корректно обрабатывает
некорректные входы. В контексте графа:

\begin{itemize}[leftmargin=1.5em]
  \item Определить \textbf{запрещённые состояния} (forbidden states)
  \item Найти пути, ведущие к ним
  \item Сгенерировать тесты, ожидающие исключение
\end{itemize}

\subsection{Property-Based Testing}

\textbf{Property-Based Testing (PBT)} --- подход, при котором вместо конкретных
тест-кейсов формулируются \textbf{инварианты} (свойства), которые должны
выполняться для любого входа.

\begin{lstlisting}[language=CommonLisp,caption={Пример property для потока эхо}]
;; Инвариант: поток-эхо всегда возвращает строку, равную входу
(for-all ((вход (gen-string)))
  (is (string= вход (запустить-поток "эхо" вход))))
\end{lstlisting}

\subsubsection{QuickCheck-подход}

Алгоритм QuickCheck:

\begin{enumerate}[leftmargin=1.5em]
  \item Сгенерировать случайный вход из заданного генератора
  \item Выполнить функцию
  \item Проверить свойство
  \item При неудаче --- shrink (минимизировать контрпример)
  \item Повторить N раз (обычно 100)
\end{enumerate}

\subsubsection{Применение к CL-потокам}

Для потоков \texttt{development\_system} можно сформулировать свойства:

\begin{center}
\begin{tabularx}{\textwidth}{lL}
  \toprule
  \textbf{Свойство} & \textbf{Формулировка} \\
  \midrule
  Типизация       & \texttt{(выполнить x)} всегда возвращает строку или выбрасывает исключение \\
  Терминация      & \texttt{(выполнить x)} завершается за конечное время (timeout) \\
  Идемпотентность & Для \texttt{эхо}: повторный вызов даёт тот же результат \\
  Детерминизм     & Для \texttt{эхо}: один вход $\to$ один выход \\
  \bottomrule
\end{tabularx}
\end{center}

\subsubsection{CL-библиотеки для PBT}

\begin{itemize}[leftmargin=1.5em]
  \item \texttt{fiveam} --- популярный testing framework с базовой поддержкой generators
  \item \texttt{cl-quickcheck} --- порт Haskell QuickCheck (заброшен, но работает)
  \item \texttt{check-it} --- современная альтернатива с shrinking
\end{itemize}

\subsection{Reachability Analysis}

\textbf{Reachability Analysis} --- анализ достижимости состояний в графе.

\subsubsection{Алгоритм}

\begin{lstlisting}[language=CommonLisp,caption={BFS для анализа достижимости}]
(defun достижимые-состояния (граф начало)
  "Возвращает множество всех достижимых состояний из НАЧАЛО"
  (let ((посещённые (make-hash-table :test #'equal))
        (очередь (list начало)))
    (loop while очередь
          for текущее = (pop очередь)
          unless (gethash текущее посещённые)
          do (setf (gethash текущее посещённые) t)
             (dolist (сосед (соседи граф текущее))
               (push сосед очередь)))
    (hash-table-keys посещённые)))
\end{lstlisting}

\subsubsection{Safety vs Liveness Properties}

\begin{center}
\begin{tabularx}{\textwidth}{lLL}
  \toprule
  \textbf{Тип} & \textbf{Формулировка} & \textbf{Пример в системе} \\
  \midrule
  Safety   & <<Плохое никогда не произойдёт>> & Поток не зависает бесконечно \\
  Liveness & <<Хорошее в итоге случится>>     & Запрос в итоге получит ответ \\
  \bottomrule
\end{tabularx}
\end{center}

\subsubsection{Нежелательные состояния в \texttt{development\_system}}

\begin{itemize}[leftmargin=1.5em]
  \item Поток выбросил необработанную ошибку (без \texttt{handler-case})
  \item Мастер завис (нет ответа > N секунд)
  \item Порождённый поток не компилируется
  \item LLM вернул невалидный код
  \item Бесконечный цикл в потоке
\end{itemize}

\subsection{TLA+ и формальная верификация}

\textbf{TLA+} (Temporal Logic of Actions) --- язык формальной спецификации,
разработанный Лесли Лэмпортом для описания конкурентных систем.

\subsubsection{Применимость к CL-потокам}

TLA+ описывает систему как:
\begin{itemize}[leftmargin=1.5em]
  \item \textbf{Состояние} --- значения всех переменных
  \item \textbf{Действие} --- предикат на паре (старое состояние, новое состояние)
  \item \textbf{Спецификация} --- формула темпоральной логики
\end{itemize}

\begin{warningbox}[title={Ограничения TLA+}]
TLA+ требует изоморфного описания системы на отдельном языке.
Для \texttt{development\_system} это означает дублирование логики.
Рекомендуется \textbf{lightweight-подход}: описание графа без
полного model checking.
\end{warningbox}

\subsection{Alloy для описания графа}

\textbf{Alloy} --- язык декларативного моделирования с автоматическим
поиском контрпримеров.

\subsubsection{Relational Model}

В Alloy всё --- отношения:

\begin{lstlisting}[language={},caption={Пример Alloy-спецификации}]
sig State {}
sig Flow { executes: State -> State }

pred valid_flow[f: Flow] {
  all s: State | some f.executes[s] implies
    f.executes[s] in State
}

run valid_flow for 5
\end{lstlisting}

\subsubsection{Применимость}

Alloy полезен для:
\begin{itemize}[leftmargin=1.5em]
  \item Поиска контрпримеров к инвариантам
  \item Проверки полноты графа (все состояния достижимы)
  \item Обнаружения тупиков (deadlocks)
\end{itemize}

\textit{Для MVP:} Alloy избыточен. Достаточно ручного анализа графа
с BFS/DFS.

\subsection{LLM как генератор тест-кейсов}

\subsubsection{Prompt Engineering для извлечения тестов}

\begin{lstlisting}[language={},caption={Промпт для генерации тестов из описания}]
Система описана следующим образом:
<описание системы>

Граф состояний:
- Начальное состояние: idle
- Целевые состояния: result_ready
- Запрещённые состояния: error, timeout

Сгенерируй тест-кейсы в формате:
ТЕСТ: <название>
ВХОД: <входные данные>
ОЖИДАНИЕ: <целевое состояние или exception>
\end{lstlisting}

\subsubsection{Проблема галлюцинаций}

LLM может сгенерировать несуществующие переходы или состояния.

\textbf{Митигация:}
\begin{enumerate}[leftmargin=1.5em]
  \item Валидация сгенерированных тестов против исходного графа
  \item Self-critique: попросить LLM проверить свой же вывод
  \item Человек проверяет первые N тестов
\end{enumerate}

% ═══════════════════════════════════════════════════════════════════
\section{Граф состояний development\_system}

\subsection{Полный граф цикла обработки сообщения}

На рисунке~\ref{fig:state-graph} представлен граф состояний для одного
цикла обработки входящего сообщения в \texttt{development\_system}.

\begin{figure}[H]
\centering
\begin{tikzpicture}[
  ->,
  >=Stealth,
  node distance=2.5cm and 3cm,
  state/.style={
    rectangle, rounded corners=6pt,
    draw, thick, minimum width=2.8cm, minimum height=1cm,
    font=\small\ttfamily, align=center
  },
  init/.style={state, fill=StateInit!20, draw=StateInit},
  target/.style={state, fill=StateTarget!20, draw=StateTarget},
  forbid/.style={state, fill=StateForbid!20, draw=StateForbid},
  normal/.style={state, fill=LightGray, draw=StateNormal},
  edge/.style={thick, font=\scriptsize},
]

% Состояния
\node[init]   (idle)     {idle\\(ожидание)};
\node[normal] (recv)     [below=of idle]     {command\_received\\(команда получена)};
\node[normal] (parse)    [below=of recv]     {command\_parsed\\(команда разобрана)};
\node[normal] (lookup)   [below left=2cm and 1.5cm of parse]  {flow\_lookup\\(поиск потока)};
\node[normal] (running)  [below=of lookup]   {flow\_running\\(поток выполняется)};
\node[target] (result)   [below=of running]  {result\_ready\\(результат готов)};
\node[normal] (respond)  [right=3.5cm of result] {response\_sent\\(ответ отправлен)};

% Запрещённые состояния
\node[forbid] (notfound) [below right=2cm and 1.5cm of parse] {flow\_not\_found\\(поток не найден)};
\node[forbid] (execerr)  [right=3.5cm of running] {execution\_error\\(ошибка выполнения)};
\node[forbid] (timeout)  [right=3.5cm of lookup]  {timeout\\(таймаут)};
\node[forbid] (llmerr)   [below=of timeout]       {llm\_error\\(ошибка LLM)};

% Рёбра
\draw[edge] (idle) -- node[right] {сообщение} (recv);
\draw[edge] (recv) -- node[right] {парсинг OK} (parse);
\draw[edge] (parse) -- node[left, pos=0.3] {/запустить} (lookup);
\draw[edge] (parse) -- node[right, pos=0.3] {поток не найден} (notfound);
\draw[edge] (lookup) -- node[left] {найден} (running);
\draw[edge] (lookup) -- node[above, pos=0.4] {таймаут поиска} (timeout);
\draw[edge] (running) -- node[left] {выполнить OK} (result);
\draw[edge] (running) -- node[above] {exception} (execerr);
\draw[edge] (running) to[bend right=20] node[left, pos=0.3] {LLM fail} (llmerr);
\draw[edge] (result) -- node[above] {отправка} (respond);
\draw[edge] (respond) to[bend right=40] node[right] {цикл} (idle);

% Легенда
\node[init, minimum width=1.5cm, minimum height=0.6cm]
  at (-3.5, -14) (leg1) {};
\node[right=0.2cm of leg1, font=\small] {Начальное};

\node[target, minimum width=1.5cm, minimum height=0.6cm]
  at (0.5, -14) (leg2) {};
\node[right=0.2cm of leg2, font=\small] {Целевое};

\node[forbid, minimum width=1.5cm, minimum height=0.6cm]
  at (4.5, -14) (leg3) {};
\node[right=0.2cm of leg3, font=\small] {Запрещённое};

\end{tikzpicture}
\caption{Граф состояний цикла обработки сообщения в \texttt{development\_system}}
\label{fig:state-graph}
\end{figure}

\subsection{Инварианты системы}

\begin{center}
\begin{tabularx}{\textwidth}{lL}
  \toprule
  \textbf{Инвариант} & \textbf{Формулировка} \\
  \midrule
  I1. Типизация & Функция \texttt{выполнить} возвращает строку или выбрасывает \texttt{condition} \\
  I2. Терминация & Любой вызов завершается за \texttt{*таймаут-потока*} секунд \\
  I3. Атомарность & Один поток не может изменить состояние другого во время выполнения \\
  I4. Идемпотентность & \texttt{(выполнить эхо x)} = \texttt{(выполнить эхо x)} для любого x \\
  I5. Целостность & После ошибки система возвращается в \texttt{idle} \\
  \bottomrule
\end{tabularx}
\end{center}

\subsection{Переходы и триггеры}

\begin{center}
\begin{tabularx}{\textwidth}{llL}
  \toprule
  \textbf{Из} & \textbf{В} & \textbf{Триггер} \\
  \midrule
  idle & command\_received & Входящее сообщение из Telegram \\
  command\_received & command\_parsed & Успешный парсинг команды \\
  command\_parsed & flow\_lookup & Команда \texttt{/запустить <имя>} \\
  command\_parsed & flow\_not\_found & Неизвестная команда \\
  flow\_lookup & flow\_running & Пакет и функция найдены \\
  flow\_lookup & timeout & Поиск превысил таймаут \\
  flow\_running & result\_ready & \texttt{выполнить} вернула строку \\
  flow\_running & execution\_error & \texttt{выполнить} выбросила condition \\
  flow\_running & llm\_error & LLM-запрос не удался \\
  result\_ready & response\_sent & Telegram API принял ответ \\
  response\_sent & idle & Цикл завершён \\
  \bottomrule
\end{tabularx}
\end{center}

% ═══════════════════════════════════════════════════════════════════
\section{Архитектура поток-граф-тестировщик}

\subsection{Обзор}

\texttt{Поток-граф-тестировщик} --- специализированный поток, который:

\begin{enumerate}[leftmargin=1.5em]
  \item Принимает описание системы в одном из трёх форматов
  \item Парсит описание и строит внутренний граф
  \item Применяет алгоритмы анализа графа
  \item Генерирует тест-кейсы
  \item Выполняет тесты через \texttt{запустить-поток}
  \item Возвращает структурированный отчёт
\end{enumerate}

\subsection{Три варианта входного формата}

\subsubsection{Вариант А: естественный язык}

\begin{lstlisting}[language={},caption={Описание на естественном языке}]
Система: Мастер принимает команду /запустить эхо задача
Ожидаемое: возвращает строку равную задаче
Нежелательное: пустой ответ, исключение в потоке

Система: Мастер принимает команду /запустить питон "1+1"
Ожидаемое: возвращает строку с результатом вычисления
Нежелательное: ошибка синтаксиса, таймаут
\end{lstlisting}

\textbf{Плюсы:} Минимальный порог входа, понятно любому.\\
\textbf{Минусы:} LLM должен интерпретировать, возможны ошибки.

\subsubsection{Вариант Б: JSON}

\begin{lstlisting}[language=JSON,caption={Граф в формате JSON}]
{
  "nodes": [
    {"id": "idle", "type": "initial"},
    {"id": "command_received", "type": "normal"},
    {"id": "flow_running", "type": "normal"},
    {"id": "result_ready", "type": "target"},
    {"id": "error", "type": "forbidden"}
  ],
  "edges": [
    {"from": "idle", "to": "command_received",
     "trigger": "message_received"},
    {"from": "command_received", "to": "flow_running",
     "trigger": "flow_found"},
    {"from": "command_received", "to": "error",
     "trigger": "flow_not_found"},
    {"from": "flow_running", "to": "result_ready",
     "trigger": "execute_ok"},
    {"from": "flow_running", "to": "error",
     "trigger": "exception"}
  ],
  "tests": {
    "happy_path": {
      "input": "/запустить эхо привет",
      "target": "result_ready"
    },
    "error_path": {
      "input": "/запустить несуществующий",
      "target": "error",
      "expect_exception": true
    }
  }
}
\end{lstlisting}

\textbf{Плюсы:} Точный, машиночитаемый, нет интерпретации.\\
\textbf{Минусы:} Громоздкий, требует знания формата.

\subsubsection{Вариант В: Lisp DSL}

\begin{lstlisting}[language=CommonLisp,caption={Граф в Lisp DSL}]
(defgraph мастер-цикл
  ;; Состояния
  (state :idle         :initial t)
  (state :cmd-received)
  (state :flow-running)
  (state :result-ready :target t)
  (state :error        :forbidden t)
  (state :timeout      :forbidden t)

  ;; Переходы
  (transition :idle -> :cmd-received
    :via "сообщение получено")
  (transition :cmd-received -> :flow-running
    :when "поток найден")
  (transition :cmd-received -> :error
    :when "поток не найден")
  (transition :flow-running -> :result-ready
    :when "выполнить вернула строку")
  (transition :flow-running -> :error
    :when "exception")
  (transition :flow-running -> :timeout
    :when "превышен таймаут")

  ;; Тесты
  (test "эхо возвращает вход"
    :input "/запустить эхо привет"
    :expect :result-ready
    :assert (string= результат "привет"))

  (test "несуществующий поток - ошибка"
    :input "/запустить несуществующий"
    :expect :error))
\end{lstlisting}

\textbf{Плюсы:} Естественный для системы на Lisp, расширяемый.\\
\textbf{Минусы:} Требует знания Lisp.

\begin{successbox}[title={Рекомендация}]
Для \texttt{development\_system} рекомендуется \textbf{Lisp DSL} как основной
формат. JSON --- для интероперабельности. Естественный язык --- для быстрых
прототипов через LLM.
\end{successbox}

\subsection{Алгоритм генерации тестов из графа}

\begin{enumerate}[leftmargin=1.5em]
  \item \textbf{Парсинг описания} $\to$ внутренний граф (adjacency list)
  \item \textbf{BFS от начального состояния} до всех target states $\to$ happy path tests
  \item \textbf{BFS до forbidden states} $\to$ negative tests (ожидаем exception)
  \item \textbf{Анализ unreachable states} $\to$ dead code / неполная спецификация
  \item \textbf{Генерация CL-кода} через порождатель $\to$ тест-потоки
\end{enumerate}

\begin{lstlisting}[language=CommonLisp,caption={Алгоритм генерации тестов}]
(defun сгенерировать-тесты (граф)
  "Генерирует тест-кейсы из графа состояний"
  (let ((начало (найти-начальное граф))
        (целевые (найти-целевые граф))
        (запрещённые (найти-запрещённые граф))
        (тесты nil))
    ;; Happy path: пути к целевым состояниям
    (dolist (цель целевые)
      (let ((путь (bfs-путь граф начало цель)))
        (when путь
          (push (make-тест :имя (format nil "happy-~A" цель)
                           :путь путь
                           :ожидание :success)
                тесты))))
    ;; Negative: пути к запрещённым состояниям
    (dolist (запр запрещённые)
      (let ((путь (bfs-путь граф начало запр)))
        (when путь
          (push (make-тест :имя (format nil "negative-~A" запр)
                           :путь путь
                           :ожидание :exception)
                тесты))))
    ;; Недостижимые состояния - предупреждение
    (let ((достижимые (достижимые-состояния граф начало)))
      (dolist (узел (все-узлы граф))
        (unless (member узел достижимые)
          (push (make-тест :имя (format nil "unreachable-~A" узел)
                           :путь nil
                           :ожидание :warning)
                тесты))))
    (nreverse тесты)))
\end{lstlisting}

\subsection{Интеграция с существующим мастером}

\begin{lstlisting}[language={},numbers=none,caption={Пример команд}]
;; Запуск с JSON-графом
/запустить граф-тестировщик граф={"nodes":[...], "edges":[...]}

;; Запуск с DSL
/запустить граф-тестировщик дсл=(defgraph мастер-цикл ...)

;; Запуск с чтением из файла
/запустить граф-тестировщик файл=AGENT.md

;; Результат:
;; ✅ тест-1: /запустить эхо "привет" -> "привет" PASS
;; ✅ тест-2: /запустить питон "1+1" -> "2" PASS
;; ❌ тест-3: /запустить несуществующий -> ERROR (expected) FAIL
;; 🟡 тест-4: состояние timeout - недостижимо в тестах
\end{lstlisting}

% ═══════════════════════════════════════════════════════════════════
\section{Реализация}

\subsection{Структуры данных}

\begin{lstlisting}[language=CommonLisp,caption={Базовые структуры данных}]
(defpackage :поток-граф-тестировщик
  (:use :cl)
  (:export #:выполнить
           #:парсить-граф #:сгенерировать-тесты
           #:запустить-тесты #:форматировать-отчёт))

(in-package :поток-граф-тестировщик)

;;; Узел графа
(defstruct узел
  (id nil :type symbol)
  (тип :normal :type (member :initial :normal :target :forbidden))
  (метаданные nil :type list))

;;; Ребро графа
(defstruct ребро
  (из nil :type symbol)
  (в nil :type symbol)
  (триггер "" :type string)
  (условие nil :type (or null function)))

;;; Граф
(defstruct граф
  (узлы nil :type list)    ; список структур узел
  (рёбра nil :type list)   ; список структур ребро
  (тесты nil :type list))  ; предопределённые тесты

;;; Тест-кейс
(defstruct тест-кейс
  (имя "" :type string)
  (путь nil :type list)         ; последовательность узлов
  (вход "" :type string)        ; входные данные
  (ожидание :success :type (member :success :exception :warning))
  (assert nil :type (or null function)))

;;; Результат теста
(defstruct результат-теста
  (тест nil :type тест-кейс)
  (статус :pending :type (member :pass :fail :skip :error :pending))
  (выход nil)
  (время-мс 0 :type integer)
  (сообщение "" :type string))
\end{lstlisting}

\subsection{Парсинг Lisp DSL}

\begin{lstlisting}[language=CommonLisp,caption={Макрос defgraph}]
(defmacro defgraph (имя &body определения)
  "Определяет граф состояний"
  `(progn
     (defparameter ,имя
       (make-граф
        :узлы (list ,@(mapcar #'обработать-state
                              (remove-if-not #'state-p определения)))
        :рёбра (list ,@(mapcar #'обработать-transition
                               (remove-if-not #'transition-p определения)))
        :тесты (list ,@(mapcar #'обработать-test
                               (remove-if-not #'test-p определения)))))
     ,имя))

(defun state-p (форма)
  (and (listp форма) (eq (first форма) 'state)))

(defun transition-p (форма)
  (and (listp форма) (eq (first форма) 'transition)))

(defun test-p (форма)
  (and (listp форма) (eq (first форма) 'test)))

(defun обработать-state (форма)
  ;; (state :имя :initial t :target nil :forbidden nil)
  (destructuring-bind (state id &key initial target forbidden) форма
    (declare (ignore state))
    `(make-узел :id ,id
                :тип ,(cond (initial :initial)
                            (target :target)
                            (forbidden :forbidden)
                            (t :normal)))))

(defun обработать-transition (форма)
  ;; (transition :из -> :в :when "условие" :via "триггер")
  (destructuring-bind (transition из arrow в &key when via) форма
    (declare (ignore transition arrow))
    `(make-ребро :из ,из :в ,в
                 :триггер ,(or via when ""))))
\end{lstlisting}

\subsection{Парсинг JSON}

\begin{lstlisting}[language=CommonLisp,caption={Парсинг JSON в граф}]
(defun парсить-json-граф (json-строка)
  "Парсит JSON-описание графа"
  (let* ((данные (yason:parse json-строка))
         (узлы-json (gethash "nodes" данные))
         (рёбра-json (gethash "edges" данные))
         (тесты-json (gethash "tests" данные)))
    (make-граф
     :узлы (mapcar
            (lambda (у)
              (make-узел
               :id (intern (string-upcase (gethash "id" у)) :keyword)
               :тип (intern (string-upcase
                            (or (gethash "type" у) "normal"))
                           :keyword)))
            узлы-json)
     :рёбра (mapcar
             (lambda (р)
               (make-ребро
                :из (intern (string-upcase (gethash "from" р)) :keyword)
                :в (intern (string-upcase (gethash "to" р)) :keyword)
                :триггер (or (gethash "trigger" р) "")))
             рёбра-json)
     :тесты (when тесты-json
              (maphash
               (lambda (имя спец)
                 (make-тест-кейс
                  :имя имя
                  :вход (gethash "input" спец)
                  :ожидание (if (gethash "expect_exception" спец)
                                :exception
                                :success)))
               тесты-json)))))
\end{lstlisting}

\subsection{Алгоритмы анализа графа}

\begin{lstlisting}[language=CommonLisp,caption={BFS и анализ достижимости}]
;;; Adjacency list из рёбер
(defun построить-adjacency (граф)
  (let ((adj (make-hash-table :test #'eq)))
    (dolist (ребро (граф-рёбра граф))
      (push (ребро-в ребро)
            (gethash (ребро-из ребро) adj)))
    adj))

;;; BFS поиск пути
(defun bfs-путь (граф из в)
  "Возвращает кратчайший путь от ИЗ до В или NIL"
  (let ((adj (построить-adjacency граф))
        (посещённые (make-hash-table :test #'eq))
        (родители (make-hash-table :test #'eq))
        (очередь (list из)))
    (setf (gethash из посещённые) t)
    (loop while очередь
          for текущий = (pop очередь)
          do (when (eq текущий в)
               (return-from bfs-путь
                 (восстановить-путь родители из в)))
             (dolist (сосед (gethash текущий adj))
               (unless (gethash сосед посещённые)
                 (setf (gethash сосед посещённые) t)
                 (setf (gethash сосед родители) текущий)
                 (push сосед очередь))))
    nil))

(defun восстановить-путь (родители из в)
  (let ((путь (list в)))
    (loop for текущий = (gethash (first путь) родители)
          while (and текущий (not (eq текущий из)))
          do (push текущий путь))
    (push из путь)
    путь))

;;; Поиск состояний по типу
(defun найти-по-типу (граф тип)
  (remove-if-not (lambda (у) (eq (узел-тип у) тип))
                 (граф-узлы граф)))

(defun найти-начальное (граф)
  (узел-id (first (найти-по-типу граф :initial))))

(defun найти-целевые (граф)
  (mapcar #'узел-id (найти-по-типу граф :target)))

(defun найти-запрещённые (граф)
  (mapcar #'узел-id (найти-по-типу граф :forbidden)))
\end{lstlisting}

\subsection{Выполнение тестов}

\begin{lstlisting}[language=CommonLisp,caption={Выполнение тест-кейсов}]
(defparameter *таймаут-теста* 10) ; секунды

(defun выполнить-тест (тест)
  "Выполняет один тест-кейс и возвращает результат"
  (let ((начало (get-internal-real-time))
        (выход nil)
        (статус :pending)
        (сообщение ""))
    (handler-case
        (progn
          ;; Парсим вход: /запустить <поток> <аргумент>
          (multiple-value-bind (поток аргумент)
              (парсить-команду-тест (тест-кейс-вход тест))
            (setf выход (запустить-поток поток аргумент))
            ;; Проверяем ожидание
            (case (тест-кейс-ожидание тест)
              (:success
               (if (stringp выход)
                   (progn
                     (setf статус :pass)
                     (setf сообщение (format nil "OK: ~A"
                                             (subseq выход 0
                                                    (min 50 (length выход))))))
                   (progn
                     (setf статус :fail)
                     (setf сообщение "Ожидалась строка"))))
              (:exception
               ;; Ожидали исключение, но не получили
               (setf статус :fail)
               (setf сообщение "Ожидалось исключение, получен результат")))))
      (error (e)
        ;; Получили исключение
        (case (тест-кейс-ожидание тест)
          (:exception
           (setf статус :pass)
           (setf сообщение (format nil "OK: exception ~A" (type-of e))))
          (otherwise
           (setf статус :error)
           (setf сообщение (format nil "ERROR: ~A" e))))))
    (let ((конец (get-internal-real-time)))
      (make-результат-теста
       :тест тест
       :статус статус
       :выход выход
       :время-мс (round (* 1000 (/ (- конец начало)
                                   internal-time-units-per-second)))
       :сообщение сообщение))))

(defun запустить-тесты (граф)
  "Выполняет все тесты из графа"
  (let ((тесты (append (граф-тесты граф)
                       (сгенерировать-тесты граф))))
    (mapcar #'выполнить-тест тесты)))
\end{lstlisting}

\subsection{Точка входа потока}

\begin{lstlisting}[language=CommonLisp,caption={Функция выполнить для потока}]
(defun выполнить (задача)
  "Главная точка входа потока-граф-тестировщика"
  (handler-case
      (let* ((формат (определить-формат задача))
             (граф (case формат
                     (:json (парсить-json-граф (извлечь-json задача)))
                     (:dsl  (eval (read-from-string
                                   (извлечь-dsl задача))))
                     (:file (парсить-файл (извлечь-путь задача)))
                     (t     (парсить-естественный задача))))
             (результаты (запустить-тесты граф)))
        (форматировать-отчёт результаты))
    (error (e)
      (format nil "❌ Ошибка граф-тестировщика: ~A" e))))

(defun определить-формат (задача)
  (cond ((search "граф={" задача) :json)
        ((search "дсл=(" задача) :dsl)
        ((search "файл=" задача) :file)
        (t :естественный)))

(defun форматировать-отчёт (результаты)
  (with-output-to-string (out)
    (format out "~%📊 Отчёт тестирования~%")
    (format out "━━━━━━━━━━━━━━━━━━━━━~%")
    (let ((pass 0) (fail 0) (skip 0) (error 0))
      (dolist (р результаты)
        (case (результат-теста-статус р)
          (:pass  (incf pass)
                  (format out "✅ ~A (~Dms)~%"
                          (тест-кейс-имя (результат-теста-тест р))
                          (результат-теста-время-мс р)))
          (:fail  (incf fail)
                  (format out "❌ ~A: ~A~%"
                          (тест-кейс-имя (результат-теста-тест р))
                          (результат-теста-сообщение р)))
          (:skip  (incf skip)
                  (format out "🟡 ~A: пропущен~%"
                          (тест-кейс-имя (результат-теста-тест р))))
          (:error (incf error)
                  (format out "💥 ~A: ~A~%"
                          (тест-кейс-имя (результат-теста-тест р))
                          (результат-теста-сообщение р)))))
      (format out "━━━━━━━━━━━━━━━━━━━━━~%")
      (format out "Итого: ✅~D ❌~D 🟡~D 💥~D~%"
              pass fail skip error))))
\end{lstlisting}

% ═══════════════════════════════════════════════════════════════════
\section{Property-Based и Contract Testing}

\subsection{Fiveam: базовое тестирование}

\begin{lstlisting}[language=CommonLisp,caption={Настройка fiveam для потоков}]
(ql:quickload :fiveam)

(defpackage :тесты-потоков
  (:use :cl :fiveam)
  (:export #:запустить-все))

(in-package :тесты-потоков)

(def-suite потоки-suite
  :description "Тесты для вычислительных потоков")

(in-suite потоки-suite)

;;; Базовые тесты
(test эхо-возвращает-вход
  "Поток эхо должен вернуть входную строку"
  (is (string= "привет"
               (запустить-поток "эхо" "привет")))
  (is (string= ""
               (запустить-поток "эхо" "")))
  (is (string= "123"
               (запустить-поток "эхо" "123"))))

(test несуществующий-поток-ошибка
  "Несуществующий поток должен выбросить ошибку"
  (signals error
    (запустить-поток "несуществующий-поток-xyz" "тест")))

(test поток-возвращает-строку
  "Любой поток должен вернуть строку или выбросить ошибку"
  (dolist (имя '("эхо" "уточнение" "исполнитель"))
    (let ((результат (запустить-поток имя "тест")))
      (is (stringp результат)
          "Поток ~A должен вернуть строку" имя))))
\end{lstlisting}

\subsection{Property-Based Testing с check-it}

\begin{lstlisting}[language=CommonLisp,caption={Property-based тесты}]
(ql:quickload :check-it)

(defpackage :тесты-свойств
  (:use :cl :check-it)
  (:export #:проверить-свойства))

(in-package :тесты-свойств)

;;; Генераторы
(defgenerator gen-safe-string ()
  "Генератор безопасных строк (без управляющих символов)"
  (let ((len (random 100)))
    (coerce (loop repeat len
                  collect (code-char (+ 32 (random 95))))
            'string)))

;;; Свойства
(defproperty эхо-идемпотентен
  "Эхо всегда возвращает то же самое при повторном вызове"
  ((вход (gen-safe-string)))
  (let ((результат-1 (запустить-поток "эхо" вход))
        (результат-2 (запустить-поток "эхо" вход)))
    (string= результат-1 результат-2)))

(defproperty эхо-identity
  "Эхо является identity-функцией"
  ((вход (gen-safe-string)))
  (string= вход (запустить-поток "эхо" вход)))

(defproperty все-потоки-возвращают-строку
  "Любой поток возвращает строку или выбрасывает ошибку"
  ((имя (gen-element '("эхо" "уточнение")))
   (вход (gen-safe-string)))
  (handler-case
      (stringp (запустить-поток имя вход))
    (error () t)))  ; исключение - тоже допустимый результат

;;; Запуск
(defun проверить-свойства ()
  (check эхо-идемпотентен :times 100)
  (check эхо-identity :times 100)
  (check все-потоки-возвращают-строку :times 50))
\end{lstlisting}

\subsection{Contract-Based Testing}

\begin{lstlisting}[language=CommonLisp,caption={Контракты для потоков}]
(defpackage :контракты
  (:use :cl)
  (:export #:defcontract #:with-contract #:проверить-контракт))

(in-package :контракты)

;;; Структура контракта
(defstruct контракт
  (имя nil :type symbol)
  (предусловие nil :type function)
  (постусловие nil :type function)
  (инвариант nil :type (or null function)))

;;; Хранилище контрактов
(defparameter *контракты* (make-hash-table :test #'equal))

;;; Макрос для определения контракта
(defmacro defcontract (имя-потока &key pre post invariant)
  "Определяет контракт для потока"
  `(setf (gethash ,имя-потока *контракты*)
         (make-контракт
          :имя (intern (string-upcase ,имя-потока))
          :предусловие (lambda (вход) ,pre)
          :постусловие (lambda (вход выход) ,post)
          :инвариант ,(when invariant
                        `(lambda () ,invariant)))))

;;; Пример контрактов
(defcontract "эхо"
  :pre (stringp вход)
  :post (and (stringp выход)
             (string= вход выход)))

(defcontract "питон"
  :pre (stringp вход)
  :post (stringp выход)
  :invariant (< (length выход) 100000))  ; защита от бесконечного вывода

;;; Обёртка с проверкой контракта
(defun с-контрактом (имя-потока вход)
  "Выполняет поток с проверкой контракта"
  (let ((контракт (gethash имя-потока *контракты*)))
    (when контракт
      ;; Проверяем предусловие
      (unless (funcall (контракт-предусловие контракт) вход)
        (error "Нарушено предусловие для ~A: ~A" имя-потока вход)))
    ;; Выполняем
    (let ((выход (запустить-поток имя-потока вход)))
      (when контракт
        ;; Проверяем постусловие
        (unless (funcall (контракт-постусловие контракт) вход выход)
          (error "Нарушено постусловие для ~A" имя-потока))
        ;; Проверяем инвариант
        (when (контракт-инвариант контракт)
          (unless (funcall (контракт-инвариант контракт))
            (error "Нарушен инвариант для ~A" имя-потока))))
      выход)))
\end{lstlisting}

\subsection{Автоматическая генерация тестов из контрактов}

\begin{lstlisting}[language=CommonLisp,caption={Генерация тестов из контрактов}]
(defun сгенерировать-тесты-из-контрактов ()
  "Генерирует тест-кейсы из всех зарегистрированных контрактов"
  (let ((тесты nil))
    (maphash
     (lambda (имя контракт)
       ;; Тест: предусловие выполняется для валидного входа
       (push (make-тест-кейс
              :имя (format nil "~A-pre-valid" имя)
              :вход "тестовая строка"
              :ожидание :success)
             тесты)
       ;; Тест: постусловие выполняется
       (push (make-тест-кейс
              :имя (format nil "~A-post-holds" имя)
              :вход "тест"
              :ожидание :success
              :assert (lambda (вход выход)
                        (funcall (контракт-постусловие контракт)
                                 вход выход)))
             тесты))
     *контракты*)
    (nreverse тесты)))
\end{lstlisting}

% ═══════════════════════════════════════════════════════════════════
\section{Мета-уровень: система тестирует себя}

\subsection{Рекурсивное тестирование}

\texttt{Граф-тестировщик} сам является потоком в системе. Это создаёт
интересную возможность: \textbf{система может тестировать себя}.

\begin{lstlisting}[language=CommonLisp,caption={Самотестирование графа-тестировщика}]
;; Граф для тестирования самого граф-тестировщика
(defgraph тестировщик-мета
  (state :idle :initial t)
  (state :parsing)
  (state :generating)
  (state :executing)
  (state :reporting :target t)
  (state :parse-error :forbidden t)
  (state :execution-error :forbidden t)

  (transition :idle -> :parsing
    :via "получен граф")
  (transition :parsing -> :generating
    :when "граф валиден")
  (transition :parsing -> :parse-error
    :when "синтаксическая ошибка")
  (transition :generating -> :executing
    :when "тесты сгенерированы")
  (transition :executing -> :reporting
    :when "все тесты выполнены")
  (transition :executing -> :execution-error
    :when "критическая ошибка")

  (test "самотест: пустой граф"
    :input "/запустить граф-тестировщик граф={}"
    :expect :parse-error)

  (test "самотест: простой граф"
    :input "/запустить граф-тестировщик граф={\"nodes\":[{\"id\":\"a\",\"type\":\"initial\"}],\"edges\":[]}"
    :expect :reporting))
\end{lstlisting}

\subsection{Иерархия тестирования}

\begin{figure}[H]
\centering
\begin{tikzpicture}[
  level/.style={
    rectangle, rounded corners=4pt, draw, thick,
    minimum width=4cm, minimum height=1cm,
    font=\small, align=center
  },
  arrow/.style={->, thick, >=Stealth},
]

\node[level, fill=AccentBlue!20] (l0) at (0, 0) {Уровень 0\\Потоки системы};
\node[level, fill=AccentTeal!20] (l1) at (0, -2.5) {Уровень 1\\Граф-тестировщик};
\node[level, fill=AccentOrange!20] (l2) at (0, -5) {Уровень 2\\Тест тестировщика};
\node[level, fill=GreenOk!20] (l3) at (0, -7.5) {Уровень 3\\Мета-тест};

\draw[arrow] (l1) -- node[right] {тестирует} (l0);
\draw[arrow] (l2) -- node[right] {тестирует} (l1);
\draw[arrow] (l3) -- node[right] {тестирует} (l2);
\draw[arrow, dashed] (l3) to[bend right=60] node[left] {рекурсия} (l3);

\end{tikzpicture}
\caption{Иерархия мета-тестирования}
\end{figure}

\subsection{Ограничения рекурсивного тестирования}

\begin{warningbox}[title={Проблема полноты}]
Система не может полностью протестировать себя --- это следует из теоремы
Гёделя о неполноте. Тест, проверяющий корректность всех тестов, не может
быть частью той же системы.
\end{warningbox}

\textbf{Практические ограничения:}

\begin{enumerate}[leftmargin=1.5em]
  \item \textbf{Бесконечная регрессия:} Мета-тест тоже нужно тестировать
  \item \textbf{Состояние образа:} Тесты могут менять глобальное состояние SBCL
  \item \textbf{Побочные эффекты:} LLM-запросы, файловая система
  \item \textbf{Детерминизм:} LLM-потоки недетерминированы
\end{enumerate}

\textbf{Митигации:}

\begin{enumerate}[leftmargin=1.5em]
  \item Ограничить глубину мета-тестирования (2--3 уровня)
  \item Изолировать тесты через \texttt{let} и \texttt{unwind-protect}
  \item Мокировать LLM-запросы для детерминизма
  \item Использовать \texttt{with-compilation-unit} для изоляции компиляции
\end{enumerate}

\subsection{Порождатель как генератор тест-потоков}

\begin{lstlisting}[language=CommonLisp,caption={Порождение тест-потока через порождатель}]
(defun породить-тест-поток (имя-теста спецификация)
  "Использует порождатель для создания специализированного тест-потока"
  (let ((описание
         (format nil
                 "Создай поток тест-~A.

                  Задача потока: выполнить тест согласно спецификации:
                  ~A

                  Поток должен:
                  1. Вызвать тестируемую функцию
                  2. Проверить результат
                  3. Вернуть 'PASS' или 'FAIL: причина'"
                 имя-теста
                 спецификация)))
    (запустить-поток "порождатель" описание)))

;; Пример использования
(породить-тест-поток
 "эхо-длинная-строка"
 "Проверить, что поток эхо корректно обрабатывает строки длиной 10000 символов")
\end{lstlisting}

% ═══════════════════════════════════════════════════════════════════
\section{Roadmap}

\subsection{MVP (1 день)}

\begin{itemize}[leftmargin=1.5em]
  \item[\checkmark] Файл \texttt{тесты.lisp} --- наполнить базовыми тест-кейсами
  \item[\checkmark] Простые assert-тесты для существующих потоков
  \item[\checkmark] Интеграция с fiveam
\end{itemize}

\begin{lstlisting}[language=CommonLisp,caption={MVP: базовые тесты в тесты.lisp}]
;; Минимальный набор тестов
(assert (stringp (запустить-поток "эхо" "тест")))
(assert (string= "привет" (запустить-поток "эхо" "привет")))
(assert (handler-case
            (progn (запустить-поток "несуществующий" "") nil)
          (error () t)))
\end{lstlisting}

\subsection{Неделя 1}

\begin{itemize}[leftmargin=1.5em]
  \item \texttt{поток-граф-тестировщик.lisp} --- принимает JSON-граф
  \item Парсинг JSON в структуру графа
  \item BFS для анализа достижимости
  \item Генерация и выполнение тестов
  \item Форматирование отчёта для Telegram
\end{itemize}

\subsection{Неделя 2}

\begin{itemize}[leftmargin=1.5em]
  \item Lisp DSL (\texttt{defgraph})
  \item Интеграция с порождателем: автогенерация тест-потоков
  \item LLM читает AGENT.md $\to$ строит граф $\to$ генерирует тесты
  \item Визуализация графа (TikZ / GraphViz экспорт)
\end{itemize}

\subsection{Месяц}

\begin{itemize}[leftmargin=1.5em]
  \item Property-based тестирование через check-it
  \item Контракты: предусловия, постусловия, инварианты
  \item Метаданные в \texttt{.мета} файлах потоков
  \item CI-подобный цикл: граф-тестировщик запускается при изменениях
  \item Статистика покрытия переходов
\end{itemize}

\subsection{Квартал}

\begin{itemize}[leftmargin=1.5em]
  \item Интеграция с TLA+ для формальной верификации критических потоков
  \item Mutation testing: автоматические мутации кода для проверки тестов
  \item Регрессионное тестирование с историей результатов
  \item Web-интерфейс для просмотра графов и результатов
\end{itemize}

% ═══════════════════════════════════════════════════════════════════
\section{Риски}

\begin{center}
\begin{tabularx}{\textwidth}{lLl}
  \toprule
  \textbf{Риск} & \textbf{Описание} & \textbf{Митигация} \\
  \midrule
  Галлюцинации LLM
    & LLM генерирует несуществующие переходы или состояния
    & Валидация против исходного графа \\
  Изоляция состояния
    & Тесты влияют на глобальное состояние SBCL
    & \texttt{unwind-protect}, fresh binding \\
  Рекурсивное тестирование
    & Бесконечная регрессия мета-тестов
    & Ограничение глубины (2--3 уровня) \\
  Неполнота спецификации
    & Граф не описывает все состояния системы
    & Анализ unreachable + warning \\
  Недетерминизм LLM
    & Результаты LLM-потоков нестабильны
    & Мокирование для тестов, retry logic \\
  Производительность
    & Много LLM-запросов при генерации тестов
    & Кэширование, batch, дешёвые модели \\
  \bottomrule
\end{tabularx}
\end{center}

\subsection{Стоимость и ресурсы}

\begin{center}
\begin{tabular}{lrl}
  \toprule
  \textbf{Операция} & \textbf{Токены} & \textbf{Примерная стоимость} \\
  \midrule
  Парсинг естественного языка в граф & $\sim$1000 & \$0.003 \\
  Генерация 10 тест-кейсов & $\sim$2000 & \$0.006 \\
  Выполнение теста с LLM-потоком & $\sim$500 & \$0.0015 \\
  Полный цикл (10 тестов) & $\sim$7000 & \$0.02 \\
  \bottomrule
\end{tabular}
\end{center}

Стоимость приемлема для персонального использования.

% ═══════════════════════════════════════════════════════════════════
\section{Заключение}

\subsection{Ключевые выводы}

\begin{enumerate}[leftmargin=1.5em]
  \item \textbf{Граф состояний --- мощный инструмент.} Формальное описание
    системы как графа позволяет автоматически генерировать тесты, анализировать
    достижимость и выявлять мёртвый код.

  \item \textbf{Три формата --- для разных нужд.} Естественный язык для
    быстрых прототипов, JSON для интероперабельности, Lisp DSL для полного
    контроля.

  \item \textbf{MBT + PBT = полное покрытие.} Model-Based Testing покрывает
    структуру (переходы графа), Property-Based Testing покрывает данные
    (инварианты).

  \item \textbf{Метациркулярность --- преимущество и сложность.} Система
    может тестировать себя, но это создаёт проблемы рекурсии и изоляции.

  \item \textbf{Контракты формализуют ожидания.} Явные предусловия и
    постусловия превращают неявные предположения в проверяемые свойства.
\end{enumerate}

\subsection{Критерий успеха}

\begin{successbox}[title={Через месяц система считается успешной, если:}]
\begin{itemize}[leftmargin=1em]
  \item Можно описать новый поток как граф и сразу получить тесты
  \item Все существующие потоки покрыты хотя бы базовыми тестами
  \item Граф-тестировщик ловит регрессии при изменениях в потоках
  \item Время на ручное тестирование сократилось на 50\%+
\end{itemize}
\end{successbox}

\subsection{Следующие шаги}

\begin{enumerate}[leftmargin=1.5em]
  \item Реализовать MVP: базовые тесты в \texttt{тесты.lisp} + fiveam
  \item Описать граф \texttt{development\_system} в Lisp DSL
  \item Создать \texttt{поток-граф-тестировщик.lisp}
  \item Запустить первый цикл автогенерации тестов
  \item Интегрировать с порождателем для автоматического обновления тестов
\end{enumerate}

\vspace{1cm}
\hrule
\vspace{4mm}
{\small\color{MidGray}
  Документ подготовлен для Артемия.
  Версия~1.0, 14~марта~2026.
  Проект: \texttt{development\_system}.
}

\end{document}
